% --- Archon physics formalization source begin ---
% archon:physics
% archon:covers PhyXMiniProblems/problem_phyx_mini_0854.lean
% archon:source-report reports/phyx_mini/problem_phyx_mini_0854.source.json

\chapter{Physics problem phyx\_mini\_0854}
\label{ch:PhyXMiniProblems_problem_phyx_mini_0854}

\paragraph{Problem source.}
\#\# Physical scenario

The electrons are fixed and cannot move.

\#\# Question

What is the electric potential energy of the proton?

\#\# Answer choices

- A: A: 1.2 \textbackslash{}times 10\textasciicircum{}\{19\}J

- B: B: -1.2 \textbackslash{}times 10\textasciicircum{}\{19\}J

- C: C: 2.2 \textbackslash{}times 10\textasciicircum{}\{19\}J

- D: D: -2.2 \textbackslash{}times 10\textasciicircum{}\{19\}J

\#\# Figure caption (auxiliary; use the image as primary evidence)

The image depicts two electrons and one proton with specific spatial relationships:

- **Electrons**: Represented by two blue circles with negative signs, positioned vertically aligned on the left side. The distance between the electrons is labeled as "0.50 nm" vertically.
- **Proton**: Shown as a small red circle with no specific sign, located to the right, horizontally aligned with the center of the electrons.
- **Distances**: 
  - The horizontal distance from the center of the electron pair to the proton is marked as "2.0 nm."
  - Both electrons are 0.50 nm away from the central dashed horizontal line.
- **Text**: Labels identifying the particles and distances:
  - "Electrons" is labeled next to the blue circles.
  - "Proton" is labeled near the red circle.
  - Measurements "0.50 nm" and "2.0 nm" correspond to the vertical and horizontal distances, respectively.
  
The layout suggests a spatial arrangement and measurement in a diagrammatic form, indicative of particle interactions at the nanoscale.

\#\# Dataset metadata

Electrostatics; Physical Model Grounding Reasoning, Spatial Relation Reasoning

\paragraph{Recorded answer/context.}
D: D: -2.2 \textbackslash{}times 10\textasciicircum{}\{19\}J

\paragraph{Figure/image path.}
/root/proposal\_for\_physic/science-mango/phyx\_mini\_run/phyx\_data/test\_image/854.png

\paragraph{Formalization target.}
create a compiling Lean file with sorry bodies at `PhyXMiniProblems/problem\_phyx\_mini\_0854.lean`. The Lean declarations must preserve the physical quantities, dimensions or dimensional roles, figure labels, governing-law hypotheses, and final relation expressed by this problem.
Use Mathlib/Physlib names found through LeanExplore where available. If a physics API is missing, introduce faithful local abstractions rather than scalar placeholder aliases.

\begin{theorem}[Physics formalization target]
\label{thm:physics:phyx_mini_0854:target}
\lean{PhyXMiniProblems.ProblemPhyXMini0854.problem_phyx_mini_0854}
\uses{def:physics:phyx-mini-0854:phyxminiproblems-problemphyxmini0854-lengthinmeters, def:physics:phyx-mini-0854:phyxminiproblems-problemphyxmini0854-coulombconstantinjoulemeterspercoulombsquared, def:physics:phyx-mini-0854:phyxminiproblems-problemphyxmini0854-energyinjoules, def:physics:phyx-mini-0854:phyxminiproblems-problemphyxmini0854-elementarychargeincoulombs, def:physics:phyx-mini-0854:phyxminiproblems-problemphyxmini0854-fixedelectronsprotonsetup, def:physics:phyx-mini-0854:phyxminiproblems-problemphyxmini0854-matchesthreeparticlescenario, def:physics:phyx-mini-0854:phyxminiproblems-problemphyxmini0854-electronsarefixed, def:physics:phyx-mini-0854:phyxminiproblems-problemphyxmini0854-matchessuppliedelectrostaticfigure, def:physics:phyx-mini-0854:phyxminiproblems-problemphyxmini0854-usessuppliedfiguregeometry, def:physics:phyx-mini-0854:phyxminiproblems-problemphyxmini0854-useselectronandprotoncharges, def:physics:phyx-mini-0854:phyxminiproblems-problemphyxmini0854-hasphysicalelectrostaticparameters, def:physics:phyx-mini-0854:phyxminiproblems-problemphyxmini0854-usesvacuumcoulombconstant, def:physics:phyx-mini-0854:phyxminiproblems-problemphyxmini0854-satisfiespointchargecoulombpotentialenergylaw, def:physics:phyx-mini-0854:phyxminiproblems-problemphyxmini0854-matchesdisplayedpotentialenergy, def:physics:phyx-mini-0854:phyxminiproblems-problemphyxmini0854-isuniquematchingdisplayedpotentialenergy}
The assigned autoformalize agent should translate this physics problem into Lean declarations in the covered file, with theorem and lemma proof bodies written as `by sorry`.
\end{theorem}
\begin{proof}
This is an autoformalization task, not a proof task. Produce statements that can later be proved by the physics prover without weakening the physical model.
\end{proof}

% --- Archon declaration topology begin ---
\section{Declaration topology}
\begin{definition}[Length Quantity]
  \label{def:physics:phyx-mini-0854:phyxminiproblems-problemphyxmini0854-lengthquantity}
  \lean{PhyXMiniProblems.ProblemPhyXMini0854.LengthQuantity}
  A nonnegative physical length, used for the three labeled separations
\end{definition}
\begin{proof}
  This is the declared model definition of Length Quantity.
\end{proof}
\begin{definition}[Signed Length Quantity]
  \label{def:physics:phyx-mini-0854:phyxminiproblems-problemphyxmini0854-signedlengthquantity}
  \lean{PhyXMiniProblems.ProblemPhyXMini0854.SignedLengthQuantity}
  A signed physical length, used for Cartesian position coordinates
\end{definition}
\begin{proof}
  This is the declared model definition of Signed Length Quantity.
\end{proof}
\begin{definition}[Signed Charge Quantity]
  \label{def:physics:phyx-mini-0854:phyxminiproblems-problemphyxmini0854-signedchargequantity}
  \lean{PhyXMiniProblems.ProblemPhyXMini0854.SignedChargeQuantity}
  A signed, unit-independent electric charge
\end{definition}
\begin{proof}
  This is the declared model definition of Signed Charge Quantity.
\end{proof}
\begin{definition}[coulomb Constant Dimension]
  \label{def:physics:phyx-mini-0854:phyxminiproblems-problemphyxmini0854-coulombconstantdimension}
  \lean{PhyXMiniProblems.ProblemPhyXMini0854.coulombConstantDimension}
  The dimension of Coulomb's constant, equivalently J m / C²
\end{definition}
\begin{proof}
  This is the declared model definition of coulomb Constant Dimension.
\end{proof}
\begin{definition}[Coulomb Constant Quantity]
  \label{def:physics:phyx-mini-0854:phyxminiproblems-problemphyxmini0854-coulombconstantquantity}
  \lean{PhyXMiniProblems.ProblemPhyXMini0854.CoulombConstantQuantity}
  \uses{def:physics:phyx-mini-0854:phyxminiproblems-problemphyxmini0854-coulombconstantdimension}
  A nonnegative, dimensionful value of Coulomb's constant
\end{definition}
\begin{proof}
  This is the declared model definition of Coulomb Constant Quantity.
\end{proof}
\begin{definition}[length Readout]
  \label{def:physics:phyx-mini-0854:phyxminiproblems-problemphyxmini0854-lengthreadout}
  \lean{PhyXMiniProblems.ProblemPhyXMini0854.lengthReadout}
  \uses{def:physics:phyx-mini-0854:phyxminiproblems-problemphyxmini0854-lengthquantity}
  Read a nonnegative physical length in a selected unit
\end{definition}
\begin{proof}
  This is the declared model definition of length Readout.
\end{proof}
\begin{definition}[signed Length Readout]
  \label{def:physics:phyx-mini-0854:phyxminiproblems-problemphyxmini0854-signedlengthreadout}
  \lean{PhyXMiniProblems.ProblemPhyXMini0854.signedLengthReadout}
  \uses{def:physics:phyx-mini-0854:phyxminiproblems-problemphyxmini0854-signedlengthquantity}
  Read a signed physical coordinate in a selected length unit
\end{definition}
\begin{proof}
  This is the declared model definition of signed Length Readout.
\end{proof}
\begin{definition}[length In Meters]
  \label{def:physics:phyx-mini-0854:phyxminiproblems-problemphyxmini0854-lengthinmeters}
  \lean{PhyXMiniProblems.ProblemPhyXMini0854.lengthInMeters}
  \uses{def:physics:phyx-mini-0854:phyxminiproblems-problemphyxmini0854-lengthquantity, def:physics:phyx-mini-0854:phyxminiproblems-problemphyxmini0854-lengthreadout}
  Read a nonnegative physical length in coherent-SI metres
\end{definition}
\begin{proof}
  This is the declared model definition of length In Meters.
\end{proof}
\begin{definition}[length In Nanometers]
  \label{def:physics:phyx-mini-0854:phyxminiproblems-problemphyxmini0854-lengthinnanometers}
  \lean{PhyXMiniProblems.ProblemPhyXMini0854.lengthInNanometers}
  \uses{def:physics:phyx-mini-0854:phyxminiproblems-problemphyxmini0854-lengthquantity, def:physics:phyx-mini-0854:phyxminiproblems-problemphyxmini0854-lengthreadout}
  Read a nonnegative physical length in nanometres
\end{definition}
\begin{proof}
  This is the declared model definition of length In Nanometers.
\end{proof}
\begin{definition}[signed Length In Meters]
  \label{def:physics:phyx-mini-0854:phyxminiproblems-problemphyxmini0854-signedlengthinmeters}
  \lean{PhyXMiniProblems.ProblemPhyXMini0854.signedLengthInMeters}
  \uses{def:physics:phyx-mini-0854:phyxminiproblems-problemphyxmini0854-signedlengthquantity, def:physics:phyx-mini-0854:phyxminiproblems-problemphyxmini0854-signedlengthreadout}
  Read a signed physical coordinate in coherent-SI metres
\end{definition}
\begin{proof}
  This is the declared model definition of signed Length In Meters.
\end{proof}
\begin{definition}[signed Charge In Coulombs]
  \label{def:physics:phyx-mini-0854:phyxminiproblems-problemphyxmini0854-signedchargeincoulombs}
  \lean{PhyXMiniProblems.ProblemPhyXMini0854.signedChargeInCoulombs}
  \uses{def:physics:phyx-mini-0854:phyxminiproblems-problemphyxmini0854-signedchargequantity}
  Read a signed physical charge in coherent-SI coulombs
\end{definition}
\begin{proof}
  This is the declared model definition of signed Charge In Coulombs.
\end{proof}
\begin{definition}[coulomb Constant In Joule Meters Per Coulomb Squared]
  \label{def:physics:phyx-mini-0854:phyxminiproblems-problemphyxmini0854-coulombconstantinjoulemeterspercoulombsquared}
  \lean{PhyXMiniProblems.ProblemPhyXMini0854.coulombConstantInJouleMetersPerCoulombSquared}
  \uses{def:physics:phyx-mini-0854:phyxminiproblems-problemphyxmini0854-coulombconstantquantity}
  Read Coulomb's constant in coherent J m / C² units
\end{definition}
\begin{proof}
  This is the declared model definition of coulomb Constant In Joule Meters Per Coulomb Squared.
\end{proof}
\begin{definition}[energy In Joules]
  \label{def:physics:phyx-mini-0854:phyxminiproblems-problemphyxmini0854-energyinjoules}
  \lean{PhyXMiniProblems.ProblemPhyXMini0854.energyInJoules}
  Read a signed physical energy in coherent-SI joules
\end{definition}
\begin{proof}
  This is the declared model definition of energy In Joules.
\end{proof}
\begin{definition}[elementary Charge In Coulombs]
  \label{def:physics:phyx-mini-0854:phyxminiproblems-problemphyxmini0854-elementarychargeincoulombs}
  \lean{PhyXMiniProblems.ProblemPhyXMini0854.elementaryChargeInCoulombs}
  Physlib's positive elementary-charge unit expressed in coulombs
\end{definition}
\begin{proof}
  This is the declared model definition of elementary Charge In Coulombs.
\end{proof}
\begin{definition}[Particle Label]
  \label{def:physics:phyx-mini-0854:phyxminiproblems-problemphyxmini0854-particlelabel}
  \lean{PhyXMiniProblems.ProblemPhyXMini0854.ParticleLabel}
  The two electrons and proton distinguished in the supplied image
\end{definition}
\begin{proof}
  This is the declared model definition of Particle Label.
\end{proof}
\begin{definition}[Particle Species]
  \label{def:physics:phyx-mini-0854:phyxminiproblems-problemphyxmini0854-particlespecies}
  \lean{PhyXMiniProblems.ProblemPhyXMini0854.ParticleSpecies}
  Particle species relevant to this three-particle electrostatic model
\end{definition}
\begin{proof}
  This is the declared model definition of Particle Species.
\end{proof}
\begin{definition}[Charge Sign]
  \label{def:physics:phyx-mini-0854:phyxminiproblems-problemphyxmini0854-chargesign}
  \lean{PhyXMiniProblems.ProblemPhyXMini0854.ChargeSign}
  Charge signs that can be printed inside a particle marker
\end{definition}
\begin{proof}
  This is the declared model definition of Charge Sign.
\end{proof}
\begin{definition}[Figure Color]
  \label{def:physics:phyx-mini-0854:phyxminiproblems-problemphyxmini0854-figurecolor}
  \lean{PhyXMiniProblems.ProblemPhyXMini0854.FigureColor}
  Marker colors visible in the primary raster
\end{definition}
\begin{proof}
  This is the declared model definition of Figure Color.
\end{proof}
\begin{definition}[Plane Point]
  \label{def:physics:phyx-mini-0854:phyxminiproblems-problemphyxmini0854-planepoint}
  \lean{PhyXMiniProblems.ProblemPhyXMini0854.PlanePoint}
  \uses{def:physics:phyx-mini-0854:phyxminiproblems-problemphyxmini0854-signedlengthquantity}
  A dimensionful point in the plane of the diagram
\end{definition}
\begin{proof}
  This is the declared model definition of Plane Point.
\end{proof}
\begin{definition}[Fixed Electrons Proton Figure]
  \label{def:physics:phyx-mini-0854:phyxminiproblems-problemphyxmini0854-fixedelectronsprotonfigure}
  \lean{PhyXMiniProblems.ProblemPhyXMini0854.FixedElectronsProtonFigure}
  \uses{def:physics:phyx-mini-0854:phyxminiproblems-problemphyxmini0854-particlelabel, def:physics:phyx-mini-0854:phyxminiproblems-problemphyxmini0854-particlespecies, def:physics:phyx-mini-0854:phyxminiproblems-problemphyxmini0854-chargesign, def:physics:phyx-mini-0854:phyxminiproblems-problemphyxmini0854-figurecolor}
  Literal presentation data read from image 854
\end{definition}
\begin{proof}
  This is the declared model definition of Fixed Electrons Proton Figure.
\end{proof}
\begin{definition}[Fixed Electrons Proton Setup]
  \label{def:physics:phyx-mini-0854:phyxminiproblems-problemphyxmini0854-fixedelectronsprotonsetup}
  \lean{PhyXMiniProblems.ProblemPhyXMini0854.FixedElectronsProtonSetup}
  \uses{def:physics:phyx-mini-0854:phyxminiproblems-problemphyxmini0854-lengthquantity, def:physics:phyx-mini-0854:phyxminiproblems-problemphyxmini0854-signedchargequantity, def:physics:phyx-mini-0854:phyxminiproblems-problemphyxmini0854-coulombconstantquantity, def:physics:phyx-mini-0854:phyxminiproblems-problemphyxmini0854-particlelabel, def:physics:phyx-mini-0854:phyxminiproblems-problemphyxmini0854-particlespecies, def:physics:phyx-mini-0854:phyxminiproblems-problemphyxmini0854-planepoint, def:physics:phyx-mini-0854:phyxminiproblems-problemphyxmini0854-fixedelectronsprotonfigure}
  The physical three-particle setup. The time argument is explicitly an SI seconds readout. It is retained so that the prose assertion that both electrons cannot move can be stated as constancy of their worldlines
\end{definition}
\begin{proof}
  This is the declared model definition of Fixed Electrons Proton Setup.
\end{proof}
\begin{definition}[Matches Three Particle Scenario]
  \label{def:physics:phyx-mini-0854:phyxminiproblems-problemphyxmini0854-matchesthreeparticlescenario}
  \lean{PhyXMiniProblems.ProblemPhyXMini0854.MatchesThreeParticleScenario}
  \uses{def:physics:phyx-mini-0854:phyxminiproblems-problemphyxmini0854-fixedelectronsprotonsetup}
  The setup contains precisely the particle species named by the problem
\end{definition}
\begin{proof}
  This is the declared model definition of Matches Three Particle Scenario.
\end{proof}
\begin{definition}[Electrons Are Fixed]
  \label{def:physics:phyx-mini-0854:phyxminiproblems-problemphyxmini0854-electronsarefixed}
  \lean{PhyXMiniProblems.ProblemPhyXMini0854.ElectronsAreFixed}
  \uses{def:physics:phyx-mini-0854:phyxminiproblems-problemphyxmini0854-fixedelectronsprotonsetup}
  The statement that the two electrons are fixed and cannot move
\end{definition}
\begin{proof}
  This is the declared model definition of Electrons Are Fixed.
\end{proof}
\begin{definition}[Matches Supplied Electrostatic Figure]
  \label{def:physics:phyx-mini-0854:phyxminiproblems-problemphyxmini0854-matchessuppliedelectrostaticfigure}
  \lean{PhyXMiniProblems.ProblemPhyXMini0854.MatchesSuppliedElectrostaticFigure}
  \uses{def:physics:phyx-mini-0854:phyxminiproblems-problemphyxmini0854-fixedelectronsprotonsetup}
  All unambiguous words, colors, signs, guides, and numbers in image 854
\end{definition}
\begin{proof}
  This is the declared model definition of Matches Supplied Electrostatic Figure.
\end{proof}
\begin{definition}[electron Midline Vertical Coordinate In Meters]
  \label{def:physics:phyx-mini-0854:phyxminiproblems-problemphyxmini0854-electronmidlineverticalcoordinateinmeters}
  \lean{PhyXMiniProblems.ProblemPhyXMini0854.electronMidlineVerticalCoordinateInMeters}
  \uses{def:physics:phyx-mini-0854:phyxminiproblems-problemphyxmini0854-signedlengthinmeters, def:physics:phyx-mini-0854:phyxminiproblems-problemphyxmini0854-fixedelectronsprotonsetup}
  Vertical coordinate of the midpoint between the two electrons, in metres
\end{definition}
\begin{proof}
  This is the declared model definition of electron Midline Vertical Coordinate In Meters.
\end{proof}
\begin{definition}[Uses Supplied Figure Geometry]
  \label{def:physics:phyx-mini-0854:phyxminiproblems-problemphyxmini0854-usessuppliedfiguregeometry}
  \lean{PhyXMiniProblems.ProblemPhyXMini0854.UsesSuppliedFigureGeometry}
  \uses{def:physics:phyx-mini-0854:phyxminiproblems-problemphyxmini0854-lengthinmeters, def:physics:phyx-mini-0854:phyxminiproblems-problemphyxmini0854-lengthinnanometers, def:physics:phyx-mini-0854:phyxminiproblems-problemphyxmini0854-signedlengthinmeters, def:physics:phyx-mini-0854:phyxminiproblems-problemphyxmini0854-fixedelectronsprotonsetup, def:physics:phyx-mini-0854:phyxminiproblems-problemphyxmini0854-electronmidlineverticalcoordinateinmeters}
  Calibration of the three printed distance labels to physical lengths and to the relative Cartesian geometry at the observation time. The two 0.50 nm labels are separate offsets from the dashed horizontal midline; they are not a claim that the full electron-to-electron separation is 0.50 nm
\end{definition}
\begin{proof}
  This is the declared model definition of Uses Supplied Figure Geometry.
\end{proof}
\begin{definition}[Uses Electron And Proton Charges]
  \label{def:physics:phyx-mini-0854:phyxminiproblems-problemphyxmini0854-useselectronandprotoncharges}
  \lean{PhyXMiniProblems.ProblemPhyXMini0854.UsesElectronAndProtonCharges}
  \uses{def:physics:phyx-mini-0854:phyxminiproblems-problemphyxmini0854-signedchargeincoulombs, def:physics:phyx-mini-0854:phyxminiproblems-problemphyxmini0854-elementarychargeincoulombs, def:physics:phyx-mini-0854:phyxminiproblems-problemphyxmini0854-fixedelectronsprotonsetup}
  The particles carry one negative or positive elementary charge as usual
\end{definition}
\begin{proof}
  This is the declared model definition of Uses Electron And Proton Charges.
\end{proof}
\begin{definition}[Has Physical Electrostatic Parameters]
  \label{def:physics:phyx-mini-0854:phyxminiproblems-problemphyxmini0854-hasphysicalelectrostaticparameters}
  \lean{PhyXMiniProblems.ProblemPhyXMini0854.HasPhysicalElectrostaticParameters}
  \uses{def:physics:phyx-mini-0854:phyxminiproblems-problemphyxmini0854-lengthinmeters, def:physics:phyx-mini-0854:phyxminiproblems-problemphyxmini0854-coulombconstantinjoulemeterspercoulombsquared, def:physics:phyx-mini-0854:phyxminiproblems-problemphyxmini0854-elementarychargeincoulombs, def:physics:phyx-mini-0854:phyxminiproblems-problemphyxmini0854-fixedelectronsprotonsetup}
  Positivity conditions for the physical parameters used in Coulomb's law
\end{definition}
\begin{proof}
  This is the declared model definition of Has Physical Electrostatic Parameters.
\end{proof}
\begin{definition}[Uses Vacuum Coulomb Constant]
  \label{def:physics:phyx-mini-0854:phyxminiproblems-problemphyxmini0854-usesvacuumcoulombconstant}
  \lean{PhyXMiniProblems.ProblemPhyXMini0854.UsesVacuumCoulombConstant}
  \uses{def:physics:phyx-mini-0854:phyxminiproblems-problemphyxmini0854-coulombconstantinjoulemeterspercoulombsquared, def:physics:phyx-mini-0854:phyxminiproblems-problemphyxmini0854-fixedelectronsprotonsetup}
  The dimensionful Coulomb constant agrees with Physlib's electromagnetic system and with the standard vacuum SI readout. This calibrates a universal constant and does not supply the requested proton energy
\end{definition}
\begin{proof}
  This is the declared model definition of Uses Vacuum Coulomb Constant.
\end{proof}
\begin{definition}[particle Separation In Meters]
  \label{def:physics:phyx-mini-0854:phyxminiproblems-problemphyxmini0854-particleseparationinmeters}
  \lean{PhyXMiniProblems.ProblemPhyXMini0854.particleSeparationInMeters}
  \uses{def:physics:phyx-mini-0854:phyxminiproblems-problemphyxmini0854-signedlengthinmeters, def:physics:phyx-mini-0854:phyxminiproblems-problemphyxmini0854-particlelabel, def:physics:phyx-mini-0854:phyxminiproblems-problemphyxmini0854-fixedelectronsprotonsetup}
  Euclidean separation of two particle centers at the observation time
\end{definition}
\begin{proof}
  This is the declared model definition of particle Separation In Meters.
\end{proof}
\begin{definition}[point Charge Coulomb Energy In Joules]
  \label{def:physics:phyx-mini-0854:phyxminiproblems-problemphyxmini0854-pointchargecoulombenergyinjoules}
  \lean{PhyXMiniProblems.ProblemPhyXMini0854.pointChargeCoulombEnergyInJoules}
  \uses{def:physics:phyx-mini-0854:phyxminiproblems-problemphyxmini0854-signedchargeincoulombs, def:physics:phyx-mini-0854:phyxminiproblems-problemphyxmini0854-coulombconstantinjoulemeterspercoulombsquared, def:physics:phyx-mini-0854:phyxminiproblems-problemphyxmini0854-particlelabel, def:physics:phyx-mini-0854:phyxminiproblems-problemphyxmini0854-fixedelectronsprotonsetup, def:physics:phyx-mini-0854:phyxminiproblems-problemphyxmini0854-particleseparationinmeters}
  One point-charge contribution k q₁ q₂ / r, read in joules
\end{definition}
\begin{proof}
  This is the declared model definition of point Charge Coulomb Energy In Joules.
\end{proof}
\begin{definition}[Satisfies Point Charge Coulomb Potential Energy Law]
  \label{def:physics:phyx-mini-0854:phyxminiproblems-problemphyxmini0854-satisfiespointchargecoulombpotentialenergylaw}
  \lean{PhyXMiniProblems.ProblemPhyXMini0854.SatisfiesPointChargeCoulombPotentialEnergyLaw}
  \uses{def:physics:phyx-mini-0854:phyxminiproblems-problemphyxmini0854-energyinjoules, def:physics:phyx-mini-0854:phyxminiproblems-problemphyxmini0854-fixedelectronsprotonsetup, def:physics:phyx-mini-0854:phyxminiproblems-problemphyxmini0854-pointchargecoulombenergyinjoules}
  The proton's potential energy, with zero at infinite separation, is the sum of its pairwise Coulomb interactions with the two fixed electrons. This is the governing law, not the numerical answer to the current question
\end{definition}
\begin{proof}
  This is the declared model definition of Satisfies Point Charge Coulomb Potential Energy Law.
\end{proof}
\begin{lemma}[electron Proton Distances from Figure]
  \label{lem:physics:phyx-mini-0854:phyxminiproblems-problemphyxmini0854-electronprotondistances-fromfigure}
  \lean{PhyXMiniProblems.ProblemPhyXMini0854.electronProtonDistances_fromFigure}
  \uses{def:physics:phyx-mini-0854:phyxminiproblems-problemphyxmini0854-lengthinmeters, def:physics:phyx-mini-0854:phyxminiproblems-problemphyxmini0854-fixedelectronsprotonsetup, def:physics:phyx-mini-0854:phyxminiproblems-problemphyxmini0854-usessuppliedfiguregeometry, def:physics:phyx-mini-0854:phyxminiproblems-problemphyxmini0854-particleseparationinmeters}
  The two electron-proton distances derived from the figure geometry
\end{lemma}
\begin{proof}
  The conclusion follows from the governing relations and figure readouts named in the statement of electron Proton Distances from Figure.
\end{proof}
\begin{lemma}[proton Potential Energy coulomb Formula]
  \label{lem:physics:phyx-mini-0854:phyxminiproblems-problemphyxmini0854-protonpotentialenergy-coulombformula}
  \lean{PhyXMiniProblems.ProblemPhyXMini0854.protonPotentialEnergy_coulombFormula}
  \uses{def:physics:phyx-mini-0854:phyxminiproblems-problemphyxmini0854-lengthinmeters, def:physics:phyx-mini-0854:phyxminiproblems-problemphyxmini0854-coulombconstantinjoulemeterspercoulombsquared, def:physics:phyx-mini-0854:phyxminiproblems-problemphyxmini0854-energyinjoules, def:physics:phyx-mini-0854:phyxminiproblems-problemphyxmini0854-elementarychargeincoulombs, def:physics:phyx-mini-0854:phyxminiproblems-problemphyxmini0854-fixedelectronsprotonsetup, def:physics:phyx-mini-0854:phyxminiproblems-problemphyxmini0854-usessuppliedfiguregeometry, def:physics:phyx-mini-0854:phyxminiproblems-problemphyxmini0854-useselectronandprotoncharges, def:physics:phyx-mini-0854:phyxminiproblems-problemphyxmini0854-satisfiespointchargecoulombpotentialenergylaw}
  Before rounding to a displayed choice, Coulomb's law and the elementary charges give this exact SI expression for the proton energy
\end{lemma}
\begin{proof}
  The conclusion follows from the governing relations and figure readouts named in the statement of proton Potential Energy coulomb Formula.
\end{proof}
\begin{definition}[Answer Choice]
  \label{def:physics:phyx-mini-0854:phyxminiproblems-problemphyxmini0854-answerchoice}
  \lean{PhyXMiniProblems.ProblemPhyXMini0854.AnswerChoice}
  Labels of the four answer choices supplied with the problem
\end{definition}
\begin{proof}
  This is the declared model definition of Answer Choice.
\end{proof}
\begin{definition}[displayed Potential Energy In Joules]
  \label{def:physics:phyx-mini-0854:phyxminiproblems-problemphyxmini0854-answerchoice-displayedpotentialenergyinjoules}
  \lean{PhyXMiniProblems.ProblemPhyXMini0854.AnswerChoice.displayedPotentialEnergyInJoules}
  \uses{def:physics:phyx-mini-0854:phyxminiproblems-problemphyxmini0854-answerchoice}
  Potential energies displayed by the four choices, in joules. The nanometre geometry and elementary charge make the physically meaningful exponent 10⁻¹⁹; the source transcription's missing exponent minus is documented in the task result
\end{definition}
\begin{proof}
  This is the declared model definition of displayed Potential Energy In Joules.
\end{proof}
\begin{definition}[displayed Energy Rounding Tolerance In Joules]
  \label{def:physics:phyx-mini-0854:phyxminiproblems-problemphyxmini0854-displayedenergyroundingtoleranceinjoules}
  \lean{PhyXMiniProblems.ProblemPhyXMini0854.displayedEnergyRoundingToleranceInJoules}
  Half a displayed 0.1 × 10⁻¹⁹ J unit
\end{definition}
\begin{proof}
  This is the declared model definition of displayed Energy Rounding Tolerance In Joules.
\end{proof}
\begin{definition}[Matches Displayed Potential Energy]
  \label{def:physics:phyx-mini-0854:phyxminiproblems-problemphyxmini0854-matchesdisplayedpotentialenergy}
  \lean{PhyXMiniProblems.ProblemPhyXMini0854.MatchesDisplayedPotentialEnergy}
  \uses{def:physics:phyx-mini-0854:phyxminiproblems-problemphyxmini0854-energyinjoules, def:physics:phyx-mini-0854:phyxminiproblems-problemphyxmini0854-fixedelectronsprotonsetup, def:physics:phyx-mini-0854:phyxminiproblems-problemphyxmini0854-answerchoice, def:physics:phyx-mini-0854:phyxminiproblems-problemphyxmini0854-displayedenergyroundingtoleranceinjoules}
  A modeled energy agrees with a choice at the displayed precision
\end{definition}
\begin{proof}
  This is the declared model definition of Matches Displayed Potential Energy.
\end{proof}
\begin{definition}[Is Unique Matching Displayed Potential Energy]
  \label{def:physics:phyx-mini-0854:phyxminiproblems-problemphyxmini0854-isuniquematchingdisplayedpotentialenergy}
  \lean{PhyXMiniProblems.ProblemPhyXMini0854.IsUniqueMatchingDisplayedPotentialEnergy}
  \uses{def:physics:phyx-mini-0854:phyxminiproblems-problemphyxmini0854-fixedelectronsprotonsetup, def:physics:phyx-mini-0854:phyxminiproblems-problemphyxmini0854-answerchoice, def:physics:phyx-mini-0854:phyxminiproblems-problemphyxmini0854-matchesdisplayedpotentialenergy}
  Exactly one displayed choice agrees with the modeled energy
\end{definition}
\begin{proof}
  This is the declared model definition of Is Unique Matching Displayed Potential Energy.
\end{proof}
\begin{definition}[recorded Dataset Answer]
  \label{def:physics:phyx-mini-0854:phyxminiproblems-problemphyxmini0854-recordeddatasetanswer}
  \lean{PhyXMiniProblems.ProblemPhyXMini0854.recordedDatasetAnswer}
  \uses{def:physics:phyx-mini-0854:phyxminiproblems-problemphyxmini0854-answerchoice}
  The source dataset's recorded answer label, retained as metadata only
\end{definition}
\begin{proof}
  This is the declared model definition of recorded Dataset Answer.
\end{proof}
% --- Archon declaration topology end ---
% --- Archon physics formalization source end ---
